\section{Automated Data Acquisition}

The dataset was collected from the Nexis Uni archive. Extracting this data was a challenge because the platform lacks a direct API and imposes strict download limits and manual navigation. To solve this, we built an automation tool using Playwright. This tool simulates a human user by navigating through search results and clicking pagination buttons automatically. To prevent the platform from blocking the session, we programmed the tool with randomized delays and varied movement patterns to mimic authentic human browsing. We used a specific boolean query to target drought impacts related to agriculture, livestock, and hydrology.

\section{Deterministic Feature Engineering}

The raw archive consisted of 300,000 JSON files containing significant "noise," such as HTML code and publisher metadata. We used a Regular Expression (Regex) pipeline to clean and structure this data. First, we separated the publication date from any event dates mentioned in the text. Second, we isolated publisher information, such as the name and headquarters of the newspaper. This was a critical step: by removing the publisher's location (e.g., "Amsterdam") from the text before it reached the model, we ensured the AI would not confuse the newspaper's office with the actual location of the drought impact (e.g., "Brabant").

\section{The Challenge of Duplicate Articles}

A major hurdle in news-based datasets is the high volume of redundant information. We identified two distinct types of duplication that require different handling strategies.

The first type is the \textbf{Syntactic Duplicate}. This occurs when the exact same article is published across multiple regional newspapers, which is common in Dutch media syndication. We address this during the preprocessing stage by using a method called MinHash with Locality Sensitive Hashing (LSH). This technique calculates the mathematical similarity between articles; if two articles overlap by more than 85\%, they are flagged as near-identical. In these cases, we keep only the earliest version to ensure the raw data is clean before the extraction begins.

The second type is the \textbf{Semantic Duplicate}. This occurs when different journalists write unique stories about the same physical event—for example, two different reports on the same crop failure in a specific province. We do not address this during preprocessing because these are not "errors" in the data, but rather different perspectives on the same event. Instead, we defer the handling of these duplicates to the post-processing phase. Depending on the modeling needs in Section B, we can either group these by their shared features (Date, Location, and Sector) to create a single event record or treat the number of mentions as a "magnitude" metric to measure how much media attention a specific impact received.

\section{Labelling procces}

The bounding variable for our labelling procces, is the classification. For our classifcaiton we aim to get atleast 30 labels per class, of the 6 classes. We also mianly care about article's that have relevance of 2/3/4. 

Keeping these functional requiremnts in mind we empoloy the following labelling procces. 

We first determine if the relevance of the label, if the label has high relevance, only then we assign it one of the 6 classes, and calculate the induced and or explicit nutscode.
We expect to have to label atleast 400 article's. We spread the labelling procces over 4 days (to not go insane lol).

Labeling Workflow: 


1. Read Headline/scan text: If not drought, mark relevance: 0 and STOP.
2. Scan for Sector: Look for the 6 keywords. Pick the primary one. Tag every sector that is explicitly mentioned. Read after.
Identify Geo-Mention: Type out the name (e.g., "Maas").
Assign Induction Type: Is it a Point, Admin, Fuzzy area, or River?
Final Check: Does this article actually prove an impact happened? If yes, it’s a relevance: 3+., i was doing multie claissfication right.




\begin{table}[h!]
\centering
\caption{Drought Impact Classification Rubric (Module A.3)}
\label{tab:sector_rubric}
\renewcommand{\arraystretch}{1.5}
\begin{tabularx}{\textwidth}{@{} l >{\RaggedRight}X >{\RaggedRight}X @{}}
\toprule
\textbf{Sector} & \textbf{Primary Indicators} & \textbf{Key Search Terms (NL)} \\ \midrule
\textbf{Agriculture} & Crop failure, reduced yields, soil moisture deficits, irrigation bans. & \textit{Oogst, gewassen, beregening, aardappelen.} \\
\textbf{Livestock} & Feed shortages, heat stress in cattle, reduced milk production. & \textit{Veehouderij, kuilvoer, weidegang, drinkbak.} \\
\textbf{Hydrological} & Low groundwater levels, lock closures, salt-water intrusion (verzilting). & \textit{Grondwater, peilbeheer, waterschap, stuw.} \\
\textbf{Energy} & Cooling water limitations for power plants, reduced hydropower output. & \textit{Koelwater, centrale, warmtelast, stroom.} \\
\textbf{Wildfire} & Forest, heath, or peat fires; heightened fire risk alerts. & \textit{Natuurbrand, heidebrand, code rood, brandweer.} \\
\textbf{Eco / Tourism} & Fish kills, blue-green algae, park closures, leisure boat restrictions. & \textit{Vissterfte, blauwalgen, recreatie, rondvaart.} \\ \bottomrule
\end{tabularx}
\end{table}

\begin{table}[h!]
\centering
\caption{Taxonomy of Spatial Induction for Toponym Resolution (Module A.2)}
\label{tab:induction_rubric}
\renewcommand{\arraystretch}{1.4}
\begin{tabularx}{\textwidth}{@{} l >{\RaggedRight}p{3cm} >{\RaggedRight}X l @{}}
\toprule
\textbf{Induction Type} & \textbf{Source Logic} & \textbf{Example Mention} & \textbf{NUTS-3 Map} \\ \midrule
\rowcolor[gray]{0.95} \textbf{Explicit: Point} & Direct 1:1 match & "Eindhoven" & NL414 \\
\textbf{Explicit: Admin} & Official Province & "Limburg" & NL421–423 \\
\rowcolor[gray]{0.95} \textbf{Vernacular} & Fuzzy/Cultural area & "De Achterhoek" & NL225 \\
\textbf{Topological} & Linear Feature & "De Waal (River)" & [List of 5+] \\ \bottomrule
\end{tabularx}
\end{table}

\begin{table}[h!]
\centering
\begin{tabular}{|l|l|p{8.5cm}|}
\hline
\textbf{Level} & \textbf{Label} & \textbf{Linguistic \& Contextual Criteria} \\ \hline
0 & Irrelevant & Metaphorical usage or topics unrelated to environmental hazards. \\ \hline
1 & Mention & Direct reference to climatic drought without associated impacts. \\ \hline
2 & Minor & Indirect impacts or preventative socio-political measures. \\ \hline
3 & Significant & Documented sectoral damage (e.g., agriculture, shipping). \\ \hline
4 & Severe & Systemic failure, national emergencies, or ecological collapse. \\ \hline
\end{tabular}
\caption{The 5-Level Drought Impact Relevance Scale.}
\end{table}



Json example to be deleted later: 

[
  {
    "id": "NEWS_2022_001",
    "meta": {
      "date": "2022-07-15",
      "source": "NOS",
      "original_filename": "raw_article_882.txt"
    },
    "text_content": "The water level in the Waal is reaching historic lows, forcing container ships to reduce loads. Farmers in the Achterhoek are also facing irrigation bans for their potato crops.",
    "labels": {
      "relevance_score": 3,
      "sectors": [
        "Agriculture",
        "Hydrological",
        "Eco / Tourism"
      ],
      "spatial_targets": [
        {
          "mention": "De Waal",
          "induction_type": "Topological",
          "nuts3_codes": ["NL224", "NL226", "NL33A"],
          "confidence_manual": 1.0
        },
        {
          "mention": "De Achterhoek",
          "induction_type": "Vernacular",
          "nuts3_codes": ["NL225"],
          "confidence_manual": 1.0 [1.0 (Certain), 0.7 (Likely), 0.5 (Guess)]
        }
      ]
    },
    "manual_notes": "Dual impact on shipping (hydro) and potatoes (agri). Achterhoek is a fuzzy region match."
  }
]





%%

% \section{Data scraping}

% We use lexus uni + playwright scarper,
% with the following logic to peform the data-scraping.


% our first thesis run, we attempted to mine all newspaper articles in Lexus
% Uni. This proved to be highly repetitive, as Nexis Uni has download limits
% and requires manually clicking through all the pages. For this purpose, we
% built a semi-automation using Playwright. This semi-automation acts as an
% autoclicker that manually clicks through pages of the website automatically.
% We also included some stochastic variation in the clicks to simulate human
% behaviour.
% In the end, the semi-automation we built turned out to be redundant be-
% cause of our premise that a smaller set of newspaper articles is already quite
% representative of the entire set of all newspaper articles. The search term on lenxu unis is drought AND (impact, agriculture, bla, bla, bla)

% [improve this section]

% \section{Data-preprocessing}
% we have like 300k article's.

% Most json's in our news-paper archive's are like this.

% We extract all feature's with regrex:
% 1. Upper date (meaning x)
% 2. Lower date (meaning y)
% 3. Uitgever
% 4. Uitgever Location.

% Afterwrds, to not consuse the model with for example uitgever location, we only keep the core article text and we normalize the text by ...

% \section{De-Duplication}

% In our article's there are two kinds of duplicate's. The first we call near identical duplicate's, where bassically the same auther, use's the exact same story, with almost the exact but at a different publisher. We will use ... (@gemeni, think of something).




% The second kind of duplicate are semenatic duplicate's, an news-article talks about a specif drought event, and impact, another write / person talks about exact same impact, but in another publisher. Two authors could talk about crop-yield reductions, in the region brababant, but in a different way/strucutre.
% For now we hope to differentie, these kind of article's, with our model feature's. If the model selects the same region, date and classifies the same type of impact. We could deduplicate it that way. Or we could use our ambiougus model, create more classifciaiotns, and potentiall enfoceres more uniquess this way. Or we could just count the mentions, and at it to somekidn of improtnce metric.








% % include data scarper
% % include data engerning (regrex stuff)

